\documentclass[11pt]{ctexart}
\usepackage[a4paper,margin=1in]{geometry}
\usepackage{longtable,booktabs}
\title{Geant4-Agent 回归测试报告（中文）}
\author{自动化评测流程}
\date{2026-03-06\_134702}
\begin{document}
\maketitle
\section{统一 Pass/Fail 总表}
\begin{longtable}{p{0.06\linewidth}p{0.08\linewidth}p{0.07\linewidth}p{0.08\linewidth}p{0.07\linewidth}p{0.07\linewidth}p{0.07\linewidth}p{0.07\linewidth}p{0.09\linewidth}p{0.06\linewidth}p{0.21\linewidth}}
\toprule
ID & 领域 & 风格 & 参数完整性 & 期望初始 & 实际初始 & 期望最终 & 实际最终 & 缺参(初->终) & 结果 & 失败原因 \\
\midrule
V2C01 & medical\_radiotherapy & precise & full & 通过 & 通过 & 通过 & 通过 & 0->0 & 通过 & - \\
V2C02 & medical\_imaging & fuzzy & full & 通过 & 通过 & 通过 & 通过 & 0->0 & 通过 & - \\
V2C03 & aerospace\_see & precise & full & 通过 & 通过 & 通过 & 通过 & 0->0 & 通过 & - \\
V2C04 & aerospace\_shielding & fuzzy & full & 通过 & 通过 & 通过 & 通过 & 0->0 & 通过 & - \\
V2C05 & semiconductor\_reliability & precise & full & 通过 & 通过 & 通过 & 通过 & 0->0 & 通过 & - \\
V2C06 & industrial\_ndt & fuzzy & full & 通过 & 通过 & 通过 & 通过 & 0->0 & 通过 & - \\
V2C07 & security\_screening & precise & full & 通过 & 通过 & 通过 & 通过 & 0->0 & 通过 & - \\
V2C08 & nuclear\_engineering & fuzzy & full & 通过 & 通过 & 通过 & 通过 & 0->0 & 通过 & - \\
V2C09 & environmental\_radiation & precise & full & 通过 & 通过 & 通过 & 通过 & 0->0 & 通过 & - \\
V2C10 & high\_energy\_physics & fuzzy & full & 通过 & 通过 & 通过 & 通过 & 0->0 & 通过 & - \\
V2C11 & education\_demo & precise & full & 通过 & 通过 & 通过 & 通过 & 0->0 & 通过 & - \\
V2C12 & oil\_logging & fuzzy & full & 通过 & 通过 & 通过 & 通过 & 0->0 & 通过 & - \\
V2C13 & marine\_radiation & precise & full & 通过 & 通过 & 通过 & 通过 & 0->0 & 通过 & - \\
V2C14 & agriculture\_irradiation & fuzzy & full & 通过 & 通过 & 通过 & 通过 & 0->0 & 通过 & - \\
V2C15 & food\_sterilization & precise & full & 通过 & 失败 & 通过 & 失败 & 1->1 & 失败 & initial\_complete\_mismatch, final\_complete\_mismatch, expected\_final\_mismatch:geometry.structure \\
V2C16 & archaeology\_scan & fuzzy & full & 通过 & 通过 & 通过 & 通过 & 0->0 & 通过 & - \\
V2C17 & geology\_assay & precise & full & 通过 & 通过 & 通过 & 通过 & 0->0 & 通过 & - \\
V2C18 & battery\_research & fuzzy & full & 通过 & 通过 & 通过 & 通过 & 0->0 & 通过 & - \\
V2C19 & automotive\_electronics & precise & full & 通过 & 通过 & 通过 & 通过 & 0->0 & 通过 & - \\
V2C20 & fusion\_diagnostics & fuzzy & full & 通过 & 通过 & 通过 & 通过 & 0->0 & 通过 & - \\
V2C21 & space\_weather & precise & full & 通过 & 通过 & 通过 & 通过 & 0->0 & 通过 & - \\
V2C22 & dosimetry\_lab & fuzzy & full & 通过 & 通过 & 通过 & 通过 & 0->0 & 通过 & - \\
V2C23 & detector\_rnd & precise & full & 通过 & 通过 & 通过 & 通过 & 0->0 & 通过 & - \\
V2C24 & accelerator\_ops & fuzzy & full & 通过 & 通过 & 通过 & 通过 & 0->0 & 通过 & - \\
V2C25 & emergency\_response & precise & full & 通过 & 通过 & 通过 & 通过 & 0->0 & 通过 & - \\
V2C26 & medical\_nuclear\_medicine & fuzzy & missing & 失败 & 失败 & 通过 & 通过 & 6->0 & 通过 & - \\
V2C27 & medical\_proton\_therapy & precise & missing & 失败 & 失败 & 通过 & 通过 & 12->0 & 失败 & expected\_final\_mismatch:geometry.structure \\
V2C28 & satellite\_payload & fuzzy & missing & 失败 & 失败 & 通过 & 通过 & 13->0 & 通过 & - \\
V2C29 & spacecraft\_habitat & precise & missing & 失败 & 失败 & 通过 & 通过 & 8->0 & 通过 & - \\
V2C30 & chip\_single\_event & fuzzy & missing & 失败 & 失败 & 通过 & 通过 & 13->0 & 通过 & - \\
V2C31 & cargo\_portal & precise & missing & 失败 & 失败 & 通过 & 通过 & 13->0 & 通过 & - \\
V2C32 & airport\_ct & fuzzy & missing & 失败 & 失败 & 通过 & 通过 & 6->0 & 通过 & - \\
V2C33 & reactor\_shielding & precise & missing & 失败 & 失败 & 通过 & 失败 & 12->1 & 失败 & final\_complete\_mismatch, expected\_final\_mismatch:geometry.structure \\
V2C34 & waste\_monitoring & fuzzy & missing & 失败 & 失败 & 通过 & 通过 & 13->0 & 通过 & - \\
V2C35 & atmospheric\_science & precise & missing & 失败 & 失败 & 通过 & 通过 & 5->0 & 通过 & - \\
V2C36 & particle\_detector\_calib & fuzzy & missing & 失败 & 失败 & 通过 & 通过 & 13->0 & 通过 & - \\
V2C37 & teaching\_lab & precise & missing & 失败 & 失败 & 通过 & 通过 & 13->0 & 通过 & - \\
V2C38 & well\_logging\_tools & fuzzy & missing & 失败 & 失败 & 通过 & 通过 & 13->0 & 通过 & - \\
V2C39 & ocean\_sensor\_hardening & precise & missing & 失败 & 失败 & 通过 & 通过 & 13->0 & 通过 & - \\
V2C40 & crop\_mutation\_study & fuzzy & missing & 失败 & 失败 & 通过 & 失败 & 13->1 & 失败 & final\_complete\_mismatch, expected\_final\_mismatch:materials.selected\_materials.0 \\
V2C41 & food\_packaging\_validation & precise & missing & 失败 & 失败 & 通过 & 通过 & 13->0 & 通过 & - \\
V2C42 & heritage\_conservation & fuzzy & missing & 失败 & 失败 & 通过 & 通过 & 13->0 & 通过 & - \\
V2C43 & mineral\_screening & precise & missing & 失败 & 失败 & 通过 & 通过 & 13->0 & 通过 & - \\
V2C44 & ev\_battery\_safety & fuzzy & missing & 失败 & 失败 & 通过 & 通过 & 13->0 & 通过 & - \\
V2C45 & autonomous\_vehicle\_radiation & precise & missing & 失败 & 失败 & 通过 & 通过 & 13->0 & 通过 & - \\
V2C46 & tokamak\_diagnostics & fuzzy & missing & 失败 & 失败 & 通过 & 通过 & 13->0 & 通过 & - \\
V2C47 & lunar\_surface\_ops & precise & missing & 失败 & 失败 & 通过 & 通过 & 13->0 & 通过 & - \\
V2C48 & hospital\_physics\_qc & fuzzy & missing & 失败 & 失败 & 通过 & 通过 & 13->0 & 通过 & - \\
V2C49 & synchrotron\_beamline & precise & missing & 失败 & 失败 & 通过 & 通过 & 13->0 & 通过 & - \\
V2C50 & radiation\_safety\_drill & fuzzy & missing & 失败 & 失败 & 通过 & 通过 & 13->0 & 通过 & - \\
\bottomrule
\end{longtable}
\section{基准对比}
\begin{itemize}
\item pass\_rate=0.920
\item initial\_alignment\_rate=0.980
\item final\_alignment\_rate=0.940
\item missing\_reduction\_rate=1.000
\item confirm\_apply\_success\_rate=1.000
\item rollback\_after\_confirm\_rate=0.000
\item delta\_pass\_rate=-0.030
\item regressions=无
\item improvements=无
\end{itemize}
\appendix
\section{附录：完整对话与内部处理链（自然语言）}
\subsection{案例 V2C01（medical\_radiotherapy）}
\textbf{对话：}
\begin{itemize}
\item 用户：????? 80 mm x 90 mm x 30 mm??? G4\_Cu?source.type=point?particle=gamma?energy=0.50 MeV?position=(0,0,-80) mm?direction=(0,0,1)?physics=QBBC?output=root?
\item 系统：配置已完成。
\end{itemize}
\textbf{内部链路（自然语言）：}
\begin{itemize}
\item 第 1 轮：动作=finalize，阶段=materials，闭环状态=True。
\item 缺失字段：无。
\item LLM 使用=True；回退=无；后端=llm\_slot\_frame+runtime\_semantic。
\end{itemize}
\subsection{案例 V2C02（medical\_imaging）}
\textbf{对话：}
\begin{itemize}
\item 用户：I need a practical medical\_radiotherapy-style setup: target as cylinder radius 33 mm half-length 52 mm made of G4\_Al, with a beam e- source at 1.25 MeV from (0,0,-110) mm heading (0,0,1), physics list FTFP\_BERT, and json output.
\item 系统：Your configuration is complete. You have set up a medical radiotherapy-style simulation with a cylindrical G4\_Al target (radius 33 mm, half-length 52 mm), an electron beam source at 1.25 MeV, the FTFP\_BERT physics list, and JSON output.
\end{itemize}
\textbf{内部链路（自然语言）：}
\begin{itemize}
\item 第 1 轮：动作=finalize，阶段=materials，闭环状态=True。
\item 缺失字段：无。
\item LLM 使用=True；回退=无；后端=llm\_slot\_frame+runtime\_semantic。
\end{itemize}
\subsection{案例 V2C03（aerospace\_see）}
\textbf{对话：}
\begin{itemize}
\item 用户：???? ?? 53 mm??? G4\_Si?source.type=isotropic?particle=proton?energy=2.00 MeV?position=(0,0,-140) mm?direction=(0,0,1)?physics=QGSP\_BERT\_HP?output=csv?
\item 系统：配置已完成。
\end{itemize}
\textbf{内部链路（自然语言）：}
\begin{itemize}
\item 第 1 轮：动作=finalize，阶段=materials，闭环状态=True。
\item 缺失字段：无。
\item LLM 使用=True；回退=无；后端=llm\_slot\_frame+runtime\_semantic。
\end{itemize}
\subsection{案例 V2C04（aerospace\_shielding）}
\textbf{对话：}
\begin{itemize}
\item 用户：I need a practical medical\_radiotherapy-style setup: target as box 140 mm x 135 mm x 60 mm made of G4\_WATER, with a point neutron source at 2.75 MeV from (0,0,-170) mm heading (0,0,1), physics list Shielding, and hdf5 output.
\item 系统：Your configuration is complete.
\end{itemize}
\textbf{内部链路（自然语言）：}
\begin{itemize}
\item 第 1 轮：动作=finalize，阶段=materials，闭环状态=True。
\item 缺失字段：无。
\item LLM 使用=True；回退=无；后端=llm\_slot\_frame+runtime\_semantic。
\end{itemize}
\subsection{案例 V2C05（semiconductor\_reliability）}
\textbf{对话：}
\begin{itemize}
\item 用户：????? ?? 57 mm ?? 88 mm??? G4\_AIR?source.type=beam?particle=gamma?energy=3.50 MeV?position=(0,0,-200) mm?direction=(0,0,1)?physics=QBBC?output=xml?
\item 系统：配置已完成。
\end{itemize}
\textbf{内部链路（自然语言）：}
\begin{itemize}
\item 第 1 轮：动作=finalize，阶段=materials，闭环状态=True。
\item 缺失字段：无。
\item LLM 使用=True；回退=无；后端=llm\_slot\_frame+runtime\_semantic。
\end{itemize}
\subsection{案例 V2C06（industrial\_ndt）}
\textbf{对话：}
\begin{itemize}
\item 用户：I need a practical medical\_radiotherapy-style setup: target as sphere radius 80 mm made of G4\_Fe, with a isotropic e- source at 4.25 MeV from (0,0,-230) mm heading (0,0,1), physics list FTFP\_BERT, and root output.
\item 系统：Your configuration is complete.
\end{itemize}
\textbf{内部链路（自然语言）：}
\begin{itemize}
\item 第 1 轮：动作=finalize，阶段=materials，闭环状态=True。
\item 缺失字段：无。
\item LLM 使用=True；回退=无；后端=llm\_slot\_frame+runtime\_semantic。
\end{itemize}
\subsection{案例 V2C07（security\_screening）}
\textbf{对话：}
\begin{itemize}
\item 用户：????? 200 mm x 180 mm x 30 mm??? G4\_Pb?source.type=point?particle=proton?energy=5.00 MeV?position=(0,0,-80) mm?direction=(0,0,1)?physics=QGSP\_BERT\_HP?output=json?
\item 系统：配置已完成。
\end{itemize}
\textbf{内部链路（自然语言）：}
\begin{itemize}
\item 第 1 轮：动作=finalize，阶段=materials，闭环状态=True。
\item 缺失字段：无。
\item LLM 使用=True；回退=无；后端=llm\_slot\_frame+runtime\_semantic。
\end{itemize}
\subsection{案例 V2C08（nuclear\_engineering）}
\textbf{对话：}
\begin{itemize}
\item 用户：I need a practical medical\_radiotherapy-style setup: target as cylinder radius 81 mm half-length 40 mm made of G4\_W, with a beam neutron source at 5.75 MeV from (0,0,-110) mm heading (0,0,1), physics list Shielding, and csv output.
\item 系统：Your configuration is complete. You have a medical radiotherapy-style setup with a cylindrical target (81 mm radius, 40 mm half-length, made of G4\_W), a neutron beam source (5.75 MeV from (0,0,-110) mm), the Shielding physics list, and CSV output.
\end{itemize}
\textbf{内部链路（自然语言）：}
\begin{itemize}
\item 第 1 轮：动作=finalize，阶段=materials，闭环状态=True。
\item 缺失字段：无。
\item LLM 使用=True；回退=无；后端=llm\_slot\_frame+runtime\_semantic。
\end{itemize}
\subsection{案例 V2C09（environmental\_radiation）}
\textbf{对话：}
\begin{itemize}
\item 用户：???? ?? 107 mm??? G4\_Cu?source.type=isotropic?particle=gamma?energy=6.50 MeV?position=(0,0,-140) mm?direction=(0,0,1)?physics=QBBC?output=hdf5?
\item 系统：好的，配置已经完成。

根据您的输入，我已为您设置了以下参数：
*   **几何**：结构类型为嵌套圆柱体，子圆柱最大半径设为 107 mm。
*   **材料**：已选定 G4\_Cu。
*   **源**：类型为各向同性点源，粒子为伽马射线，能量为 6.50 MeV。
*   **物理列表**：已选定为 QBBC。
*   **输出**：格式为 HDF5。

所有必要信息已齐备，可以生成配置文件了。
\end{itemize}
\textbf{内部链路（自然语言）：}
\begin{itemize}
\item 第 1 轮：动作=finalize，阶段=materials，闭环状态=True。
\item 缺失字段：无。
\item LLM 使用=True；回退=无；后端=llm\_slot\_frame+runtime\_semantic。
\end{itemize}
\subsection{案例 V2C10（high\_energy\_physics）}
\textbf{对话：}
\begin{itemize}
\item 用户：I need a practical medical\_radiotherapy-style setup: target as box 100 mm x 120 mm x 60 mm made of G4\_Al, with a point e- source at 7.25 MeV from (0,0,-170) mm heading (0,0,1), physics list FTFP\_BERT, and xml output.
\item 系统：Your configuration is complete.
\end{itemize}
\textbf{内部链路（自然语言）：}
\begin{itemize}
\item 第 1 轮：动作=finalize，阶段=materials，闭环状态=True。
\item 缺失字段：无。
\item LLM 使用=True；回退=无；后端=llm\_slot\_frame+runtime\_semantic。
\end{itemize}
\subsection{案例 V2C11（education\_demo）}
\textbf{对话：}
\begin{itemize}
\item 用户：????? ?? 41 mm ?? 76 mm??? G4\_Si?source.type=beam?particle=proton?energy=0.50 MeV?position=(0,0,-200) mm?direction=(0,0,1)?physics=QGSP\_BERT\_HP?output=root?
\item 系统：配置已完成。
\end{itemize}
\textbf{内部链路（自然语言）：}
\begin{itemize}
\item 第 1 轮：动作=finalize，阶段=materials，闭环状态=True。
\item 缺失字段：无。
\item LLM 使用=True；回退=无；后端=llm\_slot\_frame+runtime\_semantic。
\end{itemize}
\subsection{案例 V2C12（oil\_logging）}
\textbf{对话：}
\begin{itemize}
\item 用户：I need a practical medical\_radiotherapy-style setup: target as sphere radius 53 mm made of G4\_WATER, with a isotropic neutron source at 1.25 MeV from (0,0,-230) mm heading (0,0,1), physics list Shielding, and json output.
\item 系统：Your configuration is complete.
\end{itemize}
\textbf{内部链路（自然语言）：}
\begin{itemize}
\item 第 1 轮：动作=finalize，阶段=materials，闭环状态=True。
\item 缺失字段：无。
\item LLM 使用=True；回退=无；后端=llm\_slot\_frame+runtime\_semantic。
\end{itemize}
\subsection{案例 V2C13（marine\_radiation）}
\textbf{对话：}
\begin{itemize}
\item 用户：????? 160 mm x 165 mm x 30 mm??? G4\_AIR?source.type=point?particle=gamma?energy=2.00 MeV?position=(0,0,-80) mm?direction=(0,0,1)?physics=QBBC?output=csv?
\item 系统：配置已完成。
\end{itemize}
\textbf{内部链路（自然语言）：}
\begin{itemize}
\item 第 1 轮：动作=finalize，阶段=materials，闭环状态=True。
\item 缺失字段：无。
\item LLM 使用=True；回退=无；后端=llm\_slot\_frame+runtime\_semantic。
\end{itemize}
\subsection{案例 V2C14（agriculture\_irradiation）}
\textbf{对话：}
\begin{itemize}
\item 用户：I need a practical medical\_radiotherapy-style setup: target as cylinder radius 65 mm half-length 112 mm made of G4\_Fe, with a beam e- source at 2.75 MeV from (0,0,-110) mm heading (0,0,1), physics list FTFP\_BERT, and hdf5 output.
\item 系统：Your configuration is complete.

You have set up a medical radiotherapy-style simulation with the following confirmed components:

*   **Geometry:** A cylindrical target with a radius of 65 mm and a half-length of 112 mm.
*   **Material:** The target is made of G4\_Fe (Iron).
*   **Source:** An electron (e-) beam with an energy of 2.75 MeV, positioned at (0,0,-110) mm and directed along (0,0,1).
*   **Physics:** The FTFP\_BERT physics list is selected.
*   **Output:** Results will be saved in HDF5 format.

All necessary parameters are set. You can proceed with the simulation.
\end{itemize}
\textbf{内部链路（自然语言）：}
\begin{itemize}
\item 第 1 轮：动作=finalize，阶段=materials，闭环状态=True。
\item 缺失字段：无。
\item LLM 使用=True；回退=无；后端=llm\_slot\_frame+runtime\_semantic。
\end{itemize}
\subsection{案例 V2C15（food\_sterilization）}
\textbf{对话：}
\begin{itemize}
\item 用户：???? ?? 80 mm??? G4\_Pb?source.type=isotropic?particle=proton?energy=3.50 MeV?position=(0,0,-140) mm?direction=(0,0,1)?physics=QGSP\_BERT\_HP?output=xml?
\item 系统：好的，当前配置进度如下：

**几何**：结构类型、子体半径(child\_rmax)和根体积名已设定，还需补充子体高度(child\_hz)。
**材料**：所选材料、材料绑定和来源已同步。
**源**：类型、粒子类型和能量已设定。
**物理**：物理列表及其选择依据和来源已同步。
**输出**：格式和路径已设定。

请补充几何参数 **child\_hz** 以继续。
\end{itemize}
\textbf{内部链路（自然语言）：}
\begin{itemize}
\item 第 1 轮：动作=summarize\_progress，阶段=geometry，闭环状态=False。
\item 缺失字段：geometry.params.child\_hz。
\item LLM 使用=True；回退=无；后端=llm\_slot\_frame+runtime\_semantic。
\end{itemize}
\subsection{案例 V2C16（archaeology\_scan）}
\textbf{对话：}
\begin{itemize}
\item 用户：I need a practical medical\_radiotherapy-style setup: target as box 220 mm x 105 mm x 60 mm made of G4\_W, with a point neutron source at 4.25 MeV from (0,0,-170) mm heading (0,0,1), physics list Shielding, and root output.
\item 系统：Your configuration is complete.
\end{itemize}
\textbf{内部链路（自然语言）：}
\begin{itemize}
\item 第 1 轮：动作=finalize，阶段=materials，闭环状态=True。
\item 缺失字段：无。
\item LLM 使用=True；回退=无；后端=llm\_slot\_frame+runtime\_semantic。
\end{itemize}
\subsection{案例 V2C17（geology\_assay）}
\textbf{对话：}
\begin{itemize}
\item 用户：????? ?? 25 mm ?? 64 mm??? G4\_Cu?source.type=beam?particle=gamma?energy=5.00 MeV?position=(0,0,-200) mm?direction=(0,0,1)?physics=QBBC?output=json?
\item 系统：配置已完成。
\end{itemize}
\textbf{内部链路（自然语言）：}
\begin{itemize}
\item 第 1 轮：动作=finalize，阶段=materials，闭环状态=True。
\item 缺失字段：无。
\item LLM 使用=True；回退=无；后端=llm\_slot\_frame+runtime\_semantic。
\end{itemize}
\subsection{案例 V2C18（battery\_research）}
\textbf{对话：}
\begin{itemize}
\item 用户：I need a practical medical\_radiotherapy-style setup: target as sphere radius 107 mm made of G4\_Al, with a isotropic e- source at 5.75 MeV from (0,0,-230) mm heading (0,0,1), physics list FTFP\_BERT, and csv output.
\item 系统：Your configuration is complete. You have a medical radiotherapy-style setup with a spherical aluminum target (107 mm radius), an isotropic electron source at 5.75 MeV, the FTFP\_BERT physics list, and CSV output.
\end{itemize}
\textbf{内部链路（自然语言）：}
\begin{itemize}
\item 第 1 轮：动作=finalize，阶段=materials，闭环状态=True。
\item 缺失字段：无。
\item LLM 使用=True；回退=无；后端=llm\_slot\_frame+runtime\_semantic。
\end{itemize}
\subsection{案例 V2C19（automotive\_electronics）}
\textbf{对话：}
\begin{itemize}
\item 用户：????? 120 mm x 150 mm x 30 mm??? G4\_Si?source.type=point?particle=proton?energy=6.50 MeV?position=(0,0,-80) mm?direction=(0,0,1)?physics=QGSP\_BERT\_HP?output=hdf5?
\item 系统：配置已全部完成！所有参数都已根据你的要求设置好了。
\end{itemize}
\textbf{内部链路（自然语言）：}
\begin{itemize}
\item 第 1 轮：动作=finalize，阶段=materials，闭环状态=True。
\item 缺失字段：无。
\item LLM 使用=True；回退=无；后端=llm\_slot\_frame+runtime\_semantic。
\end{itemize}
\subsection{案例 V2C20（fusion\_diagnostics）}
\textbf{对话：}
\begin{itemize}
\item 用户：I need a practical medical\_radiotherapy-style setup: target as cylinder radius 49 mm half-length 100 mm made of G4\_WATER, with a beam neutron source at 7.25 MeV from (0,0,-110) mm heading (0,0,1), physics list Shielding, and xml output.
\item 系统：Your configuration is complete.

You have a medical radiotherapy-style setup with a cylindrical water target (radius 49 mm, half-length 100 mm). A neutron beam source at 7.25 MeV is positioned at (0,0,-110) mm, directed along the positive Z-axis. The physics list is set to Shielding, and the output format is XML.
\end{itemize}
\textbf{内部链路（自然语言）：}
\begin{itemize}
\item 第 1 轮：动作=finalize，阶段=materials，闭环状态=True。
\item 缺失字段：无。
\item LLM 使用=True；回退=无；后端=llm\_slot\_frame+runtime\_semantic。
\end{itemize}
\subsection{案例 V2C21（space\_weather）}
\textbf{对话：}
\begin{itemize}
\item 用户：???? ?? 53 mm??? G4\_AIR?source.type=isotropic?particle=gamma?energy=0.50 MeV?position=(0,0,-140) mm?direction=(0,0,1)?physics=QBBC?output=root?
\item 系统：配置已完成。
\end{itemize}
\textbf{内部链路（自然语言）：}
\begin{itemize}
\item 第 1 轮：动作=finalize，阶段=materials，闭环状态=True。
\item 缺失字段：无。
\item LLM 使用=True；回退=无；后端=llm\_slot\_frame+runtime\_semantic。
\end{itemize}
\subsection{案例 V2C22（dosimetry\_lab）}
\textbf{对话：}
\begin{itemize}
\item 用户：I need a practical medical\_radiotherapy-style setup: target as box 180 mm x 90 mm x 60 mm made of G4\_Fe, with a point e- source at 1.25 MeV from (0,0,-170) mm heading (0,0,1), physics list FTFP\_BERT, and json output.
\item 系统：Your configuration is complete. The geometry is set to a box structure with the target dimensions 180 mm x 90 mm x 60 mm. The material is G4\_Fe, and the source is a 1.25 MeV electron point source at (0,0,-170) mm directed along (0,0,1). The physics list is FTFP\_BERT with JSON output.
\end{itemize}
\textbf{内部链路（自然语言）：}
\begin{itemize}
\item 第 1 轮：动作=finalize，阶段=materials，闭环状态=True。
\item 缺失字段：无。
\item LLM 使用=True；回退=无；后端=llm\_slot\_frame+runtime\_semantic。
\end{itemize}
\subsection{案例 V2C23（detector\_rnd）}
\textbf{对话：}
\begin{itemize}
\item 用户：????? ?? 73 mm ?? 52 mm??? G4\_Pb?source.type=beam?particle=proton?energy=2.00 MeV?position=(0,0,-200) mm?direction=(0,0,1)?physics=QGSP\_BERT\_HP?output=csv?
\item 系统：配置已完成。
\end{itemize}
\textbf{内部链路（自然语言）：}
\begin{itemize}
\item 第 1 轮：动作=finalize，阶段=materials，闭环状态=True。
\item 缺失字段：无。
\item LLM 使用=True；回退=无；后端=llm\_slot\_frame+runtime\_semantic。
\end{itemize}
\subsection{案例 V2C24（accelerator\_ops）}
\textbf{对话：}
\begin{itemize}
\item 用户：I need a practical medical\_radiotherapy-style setup: target as sphere radius 80 mm made of G4\_W, with a isotropic neutron source at 2.75 MeV from (0,0,-230) mm heading (0,0,1), physics list Shielding, and hdf5 output.
\item 系统：Your configuration is complete. You have set up a medical radiotherapy-style simulation with a spherical water target (80 mm radius), an isotropic 2.75 MeV neutron source from position (0,0,-230) mm, the Shielding physics list, and HDF5 output.
\end{itemize}
\textbf{内部链路（自然语言）：}
\begin{itemize}
\item 第 1 轮：动作=finalize，阶段=materials，闭环状态=True。
\item 缺失字段：无。
\item LLM 使用=True；回退=无；后端=llm\_slot\_frame+runtime\_semantic。
\end{itemize}
\subsection{案例 V2C25（emergency\_response）}
\textbf{对话：}
\begin{itemize}
\item 用户：????? 80 mm x 135 mm x 30 mm??? G4\_Cu?source.type=point?particle=gamma?energy=3.50 MeV?position=(0,0,-80) mm?direction=(0,0,1)?physics=QBBC?output=xml?
\item 系统：配置已完成。
\end{itemize}
\textbf{内部链路（自然语言）：}
\begin{itemize}
\item 第 1 轮：动作=finalize，阶段=materials，闭环状态=True。
\item 缺失字段：无。
\item LLM 使用=True；回退=无；后端=llm\_slot\_frame+runtime\_semantic。
\end{itemize}
\subsection{案例 V2C26（medical\_nuclear\_medicine）}
\textbf{对话：}
\begin{itemize}
\item 用户：I am doing medical\_nuclear\_medicine and need a starter simulation; ask only essential missing parameters first.
\item 系统：Here's your current setup:

**Geometry:** Updated structure and module dimensions (X, Y, Z).
**Materials:** Still requires assigning materials to volumes.
**Source:** Updated position, direction, and type (point source). Still needs particle type and energy.
**Physics:** Still needs a physics list selection.
**Output:** Still needs format and file path.

What would you like to address next?
\item 用户：I need a practical medical\_radiotherapy-style setup: target as cylinder radius 33 mm half-length 88 mm made of G4\_Al, with a beam e- source at 4.25 MeV from (0,0,-110) mm heading (0,0,1), physics list FTFP\_BERT, and root output.
\item 系统：Your current settings would be changed. Please confirm if you want to update:
- Geometry structure: from `single\_box` to `single\_tubs`
- Source type: from `point` to `beam`
\item 用户：confirm
\item 系统：Your current settings would be changed. Please confirm if you want to update:
- Geometry structure: from `single\_box` to `single\_tubs`
- Source type: from `point` to `beam`
\end{itemize}
\textbf{内部链路（自然语言）：}
\begin{itemize}
\item 第 1 轮：动作=summarize\_progress，阶段=materials，闭环状态=False。
\item 缺失字段：materials.volume\_material\_map、source.particle、source.energy、physics.physics\_list、output.format、output.path。
\item LLM 使用=True；回退=无；后端=llm\_slot\_frame+runtime\_semantic。
\item 第 2 轮：动作=confirm\_overwrite，阶段=source，闭环状态=False。
\item 缺失字段：无。
\item LLM 使用=True；回退=无；后端=llm\_slot\_frame+runtime\_semantic。
\item 第 3 轮：动作=confirm\_overwrite，阶段=physics，闭环状态=True。
\item 缺失字段：无。
\item LLM 使用=True；回退=无；后端=llm\_semantic\_frame+runtime\_semantic+confirmed\_pending\_overwrite。
\end{itemize}
\subsection{案例 V2C27（medical\_proton\_therapy）}
\textbf{对话：}
\begin{itemize}
\item 用户：??? medical\_proton\_therapy ??? Geant4 ????????????????????
\item 系统：目前配置进度如下：

**几何**：需要补充结构类型和尺寸参数。  
**材料**：需要为体积指定材料。  
**源**：已确认粒子类型、来源和选择依据；还需补充源类型和能量。  
**物理**：需要选择物理列表。  
**输出**：需要指定格式和保存路径。

请继续补充上述待填项，我会随时帮你确认或调整。
\item 用户：???? ?? 107 mm??? G4\_Si?source.type=isotropic?particle=proton?energy=5.00 MeV?position=(0,0,-140) mm?direction=(0,0,1)?physics=QGSP\_BERT\_HP?output=json?
\item 系统：配置已完成。
\end{itemize}
\textbf{内部链路（自然语言）：}
\begin{itemize}
\item 第 1 轮：动作=summarize\_progress，阶段=geometry，闭环状态=False。
\item 缺失字段：geometry.structure、geometry.params.module\_x、geometry.params.module\_y、geometry.params.module\_z、materials.volume\_material\_map、source.type、source.energy、source.position、source.direction、physics.physics\_list、output.format、output.path。
\item LLM 使用=True；回退=无；后端=llm\_slot\_frame+runtime\_semantic。
\item 第 2 轮：动作=finalize，阶段=materials，闭环状态=True。
\item 缺失字段：无。
\item LLM 使用=True；回退=无；后端=llm\_slot\_frame+runtime\_semantic。
\end{itemize}
\subsection{案例 V2C28（satellite\_payload）}
\textbf{对话：}
\begin{itemize}
\item 用户：I am doing satellite\_payload and need a starter simulation; ask only essential missing parameters first.
\item 系统：To set up your satellite payload simulation, what basic shape should the payload module be, and what is its size along the x-axis?
\item 用户：I need a practical medical\_radiotherapy-style setup: target as box 140 mm x 180 mm x 60 mm made of G4\_WATER, with a point neutron source at 5.75 MeV from (0,0,-170) mm heading (0,0,1), physics list Shielding, and csv output.
\item 系统：Your simulation configuration is complete.
\end{itemize}
\textbf{内部链路（自然语言）：}
\begin{itemize}
\item 第 1 轮：动作=ask\_clarification，阶段=geometry，闭环状态=False。
\item 缺失字段：geometry.structure、geometry.params.module\_x、geometry.params.module\_y、geometry.params.module\_z、materials.volume\_material\_map、source.type、source.particle、source.energy、source.position、source.direction、physics.physics\_list、output.format、output.path。
\item LLM 使用=False；回退=E\_LLM\_FRAME\_SCHEMA\_INVALID；后端=runtime\_semantic。
\item 第 2 轮：动作=finalize，阶段=materials，闭环状态=True。
\item 缺失字段：无。
\item LLM 使用=True；回退=无；后端=llm\_slot\_frame+runtime\_semantic。
\end{itemize}
\subsection{案例 V2C29（spacecraft\_habitat）}
\textbf{对话：}
\begin{itemize}
\item 用户：??? spacecraft\_habitat ??? Geant4 ????????????????????
\item 系统：目前我们已经配置好了几何和材料部分，接下来还需要补充源、物理和输出的相关设置。请告诉我你希望使用的源类型、粒子种类以及物理过程列表。
\item 用户：????? ?? 57 mm ?? 40 mm??? G4\_AIR?source.type=beam?particle=gamma?energy=6.50 MeV?position=(0,0,-200) mm?direction=(0,0,1)?physics=QBBC?output=hdf5?
\item 系统：检测到需要覆盖已确认的设置，请确认是否应用以下修改：
- 几何结构类型：从 single\_box 改为 single\_tubs
- 材料选择：从 G4\_Al 改为 G4\_AIR
\item 用户：确认
\item 系统：检测到需要覆盖已确认的设置，请确认是否应用以下修改：
- 几何结构类型：从 single\_box 改为 single\_tubs
- 材料选择：从 G4\_Al 改为 G4\_AIR
\end{itemize}
\textbf{内部链路（自然语言）：}
\begin{itemize}
\item 第 1 轮：动作=summarize\_progress，阶段=materials，闭环状态=False。
\item 缺失字段：source.type、source.particle、source.energy、source.position、source.direction、physics.physics\_list、output.format、output.path。
\item LLM 使用=True；回退=无；后端=llm\_slot\_frame+runtime\_semantic。
\item 第 2 轮：动作=confirm\_overwrite，阶段=source，闭环状态=False。
\item 缺失字段：无。
\item LLM 使用=True；回退=无；后端=llm\_slot\_frame+runtime\_semantic。
\item 第 3 轮：动作=confirm\_overwrite，阶段=physics，闭环状态=True。
\item 缺失字段：无。
\item LLM 使用=True；回退=无；后端=llm\_semantic\_frame+runtime\_semantic+confirmed\_pending\_overwrite。
\end{itemize}
\subsection{案例 V2C30（chip\_single\_event）}
\textbf{对话：}
\begin{itemize}
\item 用户：I am doing chip\_single\_event and need a starter simulation; ask only essential missing parameters first.
\item 系统：To set up your chip\_single\_event simulation, what basic shape should the chip geometry be, and what is its size in the X direction?
\item 用户：I need a practical medical\_radiotherapy-style setup: target as sphere radius 53 mm made of G4\_Fe, with a isotropic e- source at 7.25 MeV from (0,0,-230) mm heading (0,0,1), physics list FTFP\_BERT, and xml output.
\item 系统：Your simulation is ready. The geometry is a medical radiotherapy setup with a spherical iron target (radius 53 mm). The source is an isotropic electron beam with 7.25 MeV energy, positioned at (0,0,-230) mm and directed along the positive Z-axis. The physics list is FTFP\_BERT, and the output will be in XML format.
\end{itemize}
\textbf{内部链路（自然语言）：}
\begin{itemize}
\item 第 1 轮：动作=ask\_clarification，阶段=geometry，闭环状态=False。
\item 缺失字段：geometry.structure、geometry.params.module\_x、geometry.params.module\_y、geometry.params.module\_z、materials.volume\_material\_map、source.type、source.particle、source.energy、source.position、source.direction、physics.physics\_list、output.format、output.path。
\item LLM 使用=True；回退=无；后端=llm\_slot\_frame+runtime\_semantic。
\item 第 2 轮：动作=finalize，阶段=materials，闭环状态=True。
\item 缺失字段：无。
\item LLM 使用=True；回退=无；后端=llm\_slot\_frame+runtime\_semantic。
\end{itemize}
\subsection{案例 V2C31（cargo\_portal）}
\textbf{对话：}
\begin{itemize}
\item 用户：??? cargo\_portal ??? Geant4 ????????????????????
\item 系统：您提到的“cargo\_portal”具体是什么形状呢？比如，它是一个简单的长方体，还是一个圆柱体？另外，它的尺寸大概是多少？
\item 用户：????? 200 mm x 120 mm x 30 mm??? G4\_Pb?source.type=point?particle=proton?energy=0.50 MeV?position=(0,0,-80) mm?direction=(0,0,1)?physics=QGSP\_BERT\_HP?output=root?
\item 系统：配置已完成。
\end{itemize}
\textbf{内部链路（自然语言）：}
\begin{itemize}
\item 第 1 轮：动作=ask\_clarification，阶段=geometry，闭环状态=False。
\item 缺失字段：geometry.structure、geometry.params.module\_x、geometry.params.module\_y、geometry.params.module\_z、materials.volume\_material\_map、source.type、source.particle、source.energy、source.position、source.direction、physics.physics\_list、output.format、output.path。
\item LLM 使用=True；回退=无；后端=llm\_semantic\_frame+runtime\_semantic。
\item 第 2 轮：动作=finalize，阶段=materials，闭环状态=True。
\item 缺失字段：无。
\item LLM 使用=True；回退=无；后端=llm\_slot\_frame+runtime\_semantic。
\end{itemize}
\subsection{案例 V2C32（airport\_ct）}
\textbf{对话：}
\begin{itemize}
\item 用户：I am doing airport\_ct and need a starter simulation; ask only essential missing parameters first.
\item 系统：Here's your current setup:

**Geometry:** Updated with structure and module dimensions (X, Y, Z).
**Materials:** Updated with your selected materials and their assignment to volumes.
**Physics:** Physics list configured.
**Source:** Particle type is set. Still need: source type, energy, position, and direction.
**Output:** Still need: format and file path.

Please provide the missing source and output details to proceed.
\item 用户：I need a practical medical\_radiotherapy-style setup: target as cylinder radius 81 mm half-length 76 mm made of G4\_W, with a beam neutron source at 1.25 MeV from (0,0,-110) mm heading (0,0,1), physics list Shielding, and json output.
\item 系统：Your request would overwrite two previously confirmed settings. Please confirm this change:

*   **Geometry structure:** `single\_box` → `single\_tubs`
*   **Selected materials:** `['G4\_AIR']` → `['G4\_W']`
\item 用户：confirm
\item 系统：You've already confirmed the geometry structure and selected materials. Please confirm this specific change:

*   **Geometry structure:** `single\_box` → `single\_tubs`
*   **Selected materials:** `['G4\_AIR']` → `['G4\_W']`
\end{itemize}
\textbf{内部链路（自然语言）：}
\begin{itemize}
\item 第 1 轮：动作=summarize\_progress，阶段=materials，闭环状态=False。
\item 缺失字段：source.type、source.energy、source.position、source.direction、output.format、output.path。
\item LLM 使用=True；回退=无；后端=llm\_slot\_frame+runtime\_semantic。
\item 第 2 轮：动作=confirm\_overwrite，阶段=source，闭环状态=False。
\item 缺失字段：无。
\item LLM 使用=True；回退=无；后端=llm\_slot\_frame+runtime\_semantic。
\item 第 3 轮：动作=confirm\_overwrite，阶段=physics，闭环状态=True。
\item 缺失字段：无。
\item LLM 使用=True；回退=无；后端=llm\_semantic\_frame+runtime\_semantic+confirmed\_pending\_overwrite。
\end{itemize}
\subsection{案例 V2C33（reactor\_shielding）}
\textbf{对话：}
\begin{itemize}
\item 用户：??? reactor\_shielding ??? Geant4 ????????????????????
\item 系统：目前配置进度如下：

- **几何**：待补充几何结构类型、尺寸参数（module\_x、module\_y、module\_z）。
- **材料**：待补充体积与材料的绑定关系。
- **物理**：已根据您的要求，确认并同步了物理列表及其选择来源和依据。
- **源**：待补充源类型、粒子类型、能量等信息。
- **输出**：待补充输出格式和保存路径。

请继续补充几何、材料和源的相关参数，以便完成配置。
\item 用户：???? ?? 80 mm??? G4\_Cu?source.type=isotropic?particle=gamma?energy=2.00 MeV?position=(0,0,-140) mm?direction=(0,0,1)?physics=QBBC?output=csv?
\item 系统：检测到您想要修改已确认的配置项。请确认是否应用以下更改：物理列表将从 Shielding 改为 QBBC。
\item 用户：确认
\item 系统：检测到您想要修改已确认的配置项。请确认是否应用以下更改：物理列表将从 Shielding 改为 QBBC。
\end{itemize}
\textbf{内部链路（自然语言）：}
\begin{itemize}
\item 第 1 轮：动作=summarize\_progress，阶段=geometry，闭环状态=False。
\item 缺失字段：geometry.structure、geometry.params.module\_x、geometry.params.module\_y、geometry.params.module\_z、materials.volume\_material\_map、source.type、source.particle、source.energy、source.position、source.direction、output.format、output.path。
\item LLM 使用=True；回退=无；后端=llm\_slot\_frame+runtime\_semantic。
\item 第 2 轮：动作=confirm\_overwrite，阶段=geometry，闭环状态=False。
\item 缺失字段：geometry.params.child\_hz。
\item LLM 使用=True；回退=无；后端=llm\_slot\_frame+runtime\_semantic。
\item 第 3 轮：动作=confirm\_overwrite，阶段=geometry，闭环状态=False。
\item 缺失字段：geometry.params.child\_hz。
\item LLM 使用=True；回退=无；后端=llm\_semantic\_frame+runtime\_semantic+confirmed\_pending\_overwrite。
\end{itemize}
\subsection{案例 V2C34（waste\_monitoring）}
\textbf{对话：}
\begin{itemize}
\item 用户：I am doing waste\_monitoring and need a starter simulation; ask only essential missing parameters first.
\item 系统：To set up your waste monitoring simulation, what geometry shape would you like to use, and what should its width (x-dimension) be?
\item 用户：I need a practical medical\_radiotherapy-style setup: target as box 100 mm x 165 mm x 60 mm made of G4\_Al, with a point e- source at 2.75 MeV from (0,0,-170) mm heading (0,0,1), physics list FTFP\_BERT, and hdf5 output.
\item 系统：Your configuration is complete.
\end{itemize}
\textbf{内部链路（自然语言）：}
\begin{itemize}
\item 第 1 轮：动作=ask\_clarification，阶段=geometry，闭环状态=False。
\item 缺失字段：geometry.structure、geometry.params.module\_x、geometry.params.module\_y、geometry.params.module\_z、materials.volume\_material\_map、source.type、source.particle、source.energy、source.position、source.direction、physics.physics\_list、output.format、output.path。
\item LLM 使用=True；回退=无；后端=llm\_slot\_frame+runtime\_semantic。
\item 第 2 轮：动作=finalize，阶段=materials，闭环状态=True。
\item 缺失字段：无。
\item LLM 使用=True；回退=无；后端=llm\_slot\_frame+runtime\_semantic。
\end{itemize}
\subsection{案例 V2C35（atmospheric\_science）}
\textbf{对话：}
\begin{itemize}
\item 用户：??? atmospheric\_science ??? Geant4 ????????????????????
\item 系统：目前配置进度如下：

- **几何**：已更新结构类型和部分尺寸参数。
- **材料**：还需设置体积与材料的对应关系。
- **源**：已设定能量、位置和方向，待选择粒子类型。
- **物理**：需补充物理列表。
- **输出**：输出格式和保存路径尚未指定。

请继续补充剩余信息，完成后我会为你生成完整的配置文件。
\item 用户：????? ?? 41 mm ?? 112 mm??? G4\_Si?source.type=beam?particle=proton?energy=3.50 MeV?position=(0,0,-200) mm?direction=(0,0,1)?physics=QGSP\_BERT\_HP?output=xml?
\item 系统：检测到需要覆盖已确认的配置项，请确认是否应用以下修改：

- 几何结构类型：从 single\_box 改为 single\_tubs
- 源类型：从 point 改为 beam
\item 用户：确认
\item 系统：检测到需要覆盖已确认的配置项，请确认是否应用以下修改：

- 几何结构类型：从 single\_box 改为 single\_tubs
- 源类型：从 point 改为 beam
\end{itemize}
\textbf{内部链路（自然语言）：}
\begin{itemize}
\item 第 1 轮：动作=summarize\_progress，阶段=materials，闭环状态=False。
\item 缺失字段：materials.volume\_material\_map、source.particle、physics.physics\_list、output.format、output.path。
\item LLM 使用=True；回退=无；后端=llm\_slot\_frame+runtime\_semantic。
\item 第 2 轮：动作=confirm\_overwrite，阶段=source，闭环状态=False。
\item 缺失字段：无。
\item LLM 使用=True；回退=无；后端=llm\_slot\_frame+runtime\_semantic。
\item 第 3 轮：动作=confirm\_overwrite，阶段=physics，闭环状态=True。
\item 缺失字段：无。
\item LLM 使用=True；回退=无；后端=llm\_semantic\_frame+runtime\_semantic+confirmed\_pending\_overwrite。
\end{itemize}
\subsection{案例 V2C36（particle\_detector\_calib）}
\textbf{对话：}
\begin{itemize}
\item 用户：I am doing particle\_detector\_calib and need a starter simulation; ask only essential missing parameters first.
\item 系统：To help set up your particle detector calibration simulation, could you specify the basic shape of your detector geometry, such as a box or a tube?
\item 用户：I need a practical medical\_radiotherapy-style setup: target as sphere radius 107 mm made of G4\_WATER, with a isotropic neutron source at 4.25 MeV from (0,0,-230) mm heading (0,0,1), physics list Shielding, and root output.
\item 系统：Your particle detector calibration simulation is ready.

I've configured a medical radiotherapy-style setup with a spherical water target (107 mm radius) and an isotropic 4.25 MeV neutron source positioned at (0,0,-230) mm, directed along the positive Z-axis. The physics list is set to Shielding, and the output format is ROOT.

All essential parameters are now set. You can proceed with the simulation.
\end{itemize}
\textbf{内部链路（自然语言）：}
\begin{itemize}
\item 第 1 轮：动作=ask\_clarification，阶段=geometry，闭环状态=False。
\item 缺失字段：geometry.structure、geometry.params.module\_x、geometry.params.module\_y、geometry.params.module\_z、materials.volume\_material\_map、source.type、source.particle、source.energy、source.position、source.direction、physics.physics\_list、output.format、output.path。
\item LLM 使用=True；回退=无；后端=llm\_slot\_frame+runtime\_semantic。
\item 第 2 轮：动作=finalize，阶段=materials，闭环状态=True。
\item 缺失字段：无。
\item LLM 使用=True；回退=无；后端=llm\_slot\_frame+runtime\_semantic。
\end{itemize}
\subsection{案例 V2C37（teaching\_lab）}
\textbf{对话：}
\begin{itemize}
\item 用户：??? teaching\_lab ??? Geant4 ????????????????????
\item 系统：你提到的教学实验室几何结构具体是什么类型呢？比如是一个简单的盒子形状，还是圆柱体？另外，这个结构的横向尺寸大概是多少？
\item 用户：????? 160 mm x 105 mm x 30 mm??? G4\_AIR?source.type=point?particle=gamma?energy=5.00 MeV?position=(0,0,-80) mm?direction=(0,0,1)?physics=QBBC?output=json?
\item 系统：配置已完成。
\end{itemize}
\textbf{内部链路（自然语言）：}
\begin{itemize}
\item 第 1 轮：动作=ask\_clarification，阶段=geometry，闭环状态=False。
\item 缺失字段：geometry.structure、geometry.params.module\_x、geometry.params.module\_y、geometry.params.module\_z、materials.volume\_material\_map、source.type、source.particle、source.energy、source.position、source.direction、physics.physics\_list、output.format、output.path。
\item LLM 使用=True；回退=无；后端=llm\_semantic\_frame+runtime\_semantic。
\item 第 2 轮：动作=finalize，阶段=materials，闭环状态=True。
\item 缺失字段：无。
\item LLM 使用=True；回退=无；后端=llm\_slot\_frame+runtime\_semantic。
\end{itemize}
\subsection{案例 V2C38（well\_logging\_tools）}
\textbf{对话：}
\begin{itemize}
\item 用户：I am doing well\_logging\_tools and need a starter simulation; ask only essential missing parameters first.
\item 系统：To set up your well logging simulation, what geometry shape would you like to use — for example, a simple box or cylinder?
\item 用户：I need a practical medical\_radiotherapy-style setup: target as cylinder radius 65 mm half-length 64 mm made of G4\_Fe, with a beam e- source at 5.75 MeV from (0,0,-110) mm heading (0,0,1), physics list FTFP\_BERT, and csv output.
\item 系统：Your simulation configuration is complete.

You have set up a medical radiotherapy-style geometry with a cylindrical iron target. The electron beam source, physics list, and CSV output are all configured as you specified.
\end{itemize}
\textbf{内部链路（自然语言）：}
\begin{itemize}
\item 第 1 轮：动作=ask\_clarification，阶段=geometry，闭环状态=False。
\item 缺失字段：geometry.structure、geometry.params.module\_x、geometry.params.module\_y、geometry.params.module\_z、materials.volume\_material\_map、source.type、source.particle、source.energy、source.position、source.direction、physics.physics\_list、output.format、output.path。
\item LLM 使用=False；回退=E\_LLM\_FRAME\_SCHEMA\_INVALID；后端=runtime\_semantic。
\item 第 2 轮：动作=finalize，阶段=materials，闭环状态=True。
\item 缺失字段：无。
\item LLM 使用=True；回退=无；后端=llm\_slot\_frame+runtime\_semantic。
\end{itemize}
\subsection{案例 V2C39（ocean\_sensor\_hardening）}
\textbf{对话：}
\begin{itemize}
\item 用户：??? ocean\_sensor\_hardening ??? Geant4 ????????????????????
\item 系统：你提到的 ocean\_sensor\_hardening 在 Geant4 中建模，需要先确定它的基本形状，比如是一个方盒还是圆柱体？另外，这个模块在 X 方向的具体尺寸是多少呢？
\item 用户：???? ?? 53 mm??? G4\_Pb?source.type=isotropic?particle=proton?energy=6.50 MeV?position=(0,0,-140) mm?direction=(0,0,1)?physics=QGSP\_BERT\_HP?output=hdf5?
\item 系统：好的，你的 Geant4 配置已经完成。

根据你的输入，我已经为你设置了以下参数：
*   **几何结构**：圆柱体，半径 53 mm。
*   **材料**：G4\_Pb (铅)。
*   **源**：各向同性质子源，能量 6.50 MeV，起始位置 (0,0,-140) mm，方向 (0,0,1)。
*   **物理列表**：QGSP\_BERT\_HP。
*   **输出格式**：HDF5。

所有必要信息都已齐备，可以生成相应的配置文件了。
\end{itemize}
\textbf{内部链路（自然语言）：}
\begin{itemize}
\item 第 1 轮：动作=ask\_clarification，阶段=geometry，闭环状态=False。
\item 缺失字段：geometry.structure、geometry.params.module\_x、geometry.params.module\_y、geometry.params.module\_z、materials.volume\_material\_map、source.type、source.particle、source.energy、source.position、source.direction、physics.physics\_list、output.format、output.path。
\item LLM 使用=True；回退=无；后端=llm\_semantic\_frame+runtime\_semantic。
\item 第 2 轮：动作=finalize，阶段=materials，闭环状态=True。
\item 缺失字段：无。
\item LLM 使用=True；回退=无；后端=llm\_slot\_frame+runtime\_semantic。
\end{itemize}
\subsection{案例 V2C40（crop\_mutation\_study）}
\textbf{对话：}
\begin{itemize}
\item 用户：I am doing crop\_mutation\_study and need a starter simulation; ask only essential missing parameters first.
\item 系统：To set up your crop mutation study simulation, could you specify the basic shape of your experimental volume—for example, a simple box or cylinder?
\item 用户：I need a practical medical\_radiotherapy-style setup: target as box 220 mm x 150 mm x 60 mm made of G4\_W, with a point neutron source at 7.25 MeV from (0,0,-170) mm heading (0,0,1), physics list Shielding, and xml output.
\item 系统：Here's a summary of your setup so far:

**Geometry:** Updated with a box structure and its X (220 mm) and Y (150 mm) dimensions.
**Materials:** Still requires you to assign the material G4\_W to your target volume.
**Source:** Updated to a point neutron source with an energy of 7.25 MeV.
**Physics:** Updated to the Shielding physics list.
**Output:** Updated to XML format with a specified path.

To proceed, please confirm the material binding for your target volume.
\end{itemize}
\textbf{内部链路（自然语言）：}
\begin{itemize}
\item 第 1 轮：动作=ask\_clarification，阶段=geometry，闭环状态=False。
\item 缺失字段：geometry.structure、geometry.params.module\_x、geometry.params.module\_y、geometry.params.module\_z、materials.volume\_material\_map、source.type、source.particle、source.energy、source.position、source.direction、physics.physics\_list、output.format、output.path。
\item LLM 使用=True；回退=无；后端=llm\_slot\_frame+runtime\_semantic。
\item 第 2 轮：动作=summarize\_progress，阶段=materials，闭环状态=False。
\item 缺失字段：materials.volume\_material\_map。
\item LLM 使用=True；回退=无；后端=llm\_slot\_frame+runtime\_semantic。
\end{itemize}
\subsection{案例 V2C41（food\_packaging\_validation）}
\textbf{对话：}
\begin{itemize}
\item 用户：??? food\_packaging\_validation ??? Geant4 ????????????????????
\item 系统：你提到食品包装验证，能具体说说包装的几何形状吗？比如，它是一个简单的盒子还是圆柱体？另外，这个包装模块在X方向上的尺寸大概是多少？
\item 用户：????? ?? 25 mm ?? 100 mm??? G4\_Cu?source.type=beam?particle=gamma?energy=0.50 MeV?position=(0,0,-200) mm?direction=(0,0,1)?physics=QBBC?output=root?
\item 系统：配置已完成。
\end{itemize}
\textbf{内部链路（自然语言）：}
\begin{itemize}
\item 第 1 轮：动作=ask\_clarification，阶段=geometry，闭环状态=False。
\item 缺失字段：geometry.structure、geometry.params.module\_x、geometry.params.module\_y、geometry.params.module\_z、materials.volume\_material\_map、source.type、source.particle、source.energy、source.position、source.direction、physics.physics\_list、output.format、output.path。
\item LLM 使用=True；回退=无；后端=llm\_semantic\_frame+runtime\_semantic。
\item 第 2 轮：动作=finalize，阶段=materials，闭环状态=True。
\item 缺失字段：无。
\item LLM 使用=True；回退=无；后端=llm\_slot\_frame+runtime\_semantic。
\end{itemize}
\subsection{案例 V2C42（heritage\_conservation）}
\textbf{对话：}
\begin{itemize}
\item 用户：I am doing heritage\_conservation and need a starter simulation; ask only essential missing parameters first.
\item 系统：To set up your heritage conservation simulation, could you specify the basic shape of the object you want to model, such as a simple box or cylinder?
\item 用户：I need a practical medical\_radiotherapy-style setup: target as sphere radius 80 mm made of G4\_Al, with a isotropic e- source at 1.25 MeV from (0,0,-230) mm heading (0,0,1), physics list FTFP\_BERT, and json output.
\item 系统：Your simulation configuration is complete.
\end{itemize}
\textbf{内部链路（自然语言）：}
\begin{itemize}
\item 第 1 轮：动作=ask\_clarification，阶段=geometry，闭环状态=False。
\item 缺失字段：geometry.structure、geometry.params.module\_x、geometry.params.module\_y、geometry.params.module\_z、materials.volume\_material\_map、source.type、source.particle、source.energy、source.position、source.direction、physics.physics\_list、output.format、output.path。
\item LLM 使用=True；回退=无；后端=llm\_slot\_frame+runtime\_semantic。
\item 第 2 轮：动作=finalize，阶段=materials，闭环状态=True。
\item 缺失字段：无。
\item LLM 使用=True；回退=无；后端=llm\_slot\_frame+runtime\_semantic。
\end{itemize}
\subsection{案例 V2C43（mineral\_screening）}
\textbf{对话：}
\begin{itemize}
\item 用户：??? mineral\_screening ??? Geant4 ????????????????????
\item 系统：你希望模拟的几何结构是长方体还是圆柱体？如果是长方体，它的长度（X方向）是多少？
\item 用户：????? 120 mm x 90 mm x 30 mm??? G4\_Si?source.type=point?particle=proton?energy=2.00 MeV?position=(0,0,-80) mm?direction=(0,0,1)?physics=QGSP\_BERT\_HP?output=csv?
\item 系统：配置已完成。
\end{itemize}
\textbf{内部链路（自然语言）：}
\begin{itemize}
\item 第 1 轮：动作=ask\_clarification，阶段=geometry，闭环状态=False。
\item 缺失字段：geometry.structure、geometry.params.module\_x、geometry.params.module\_y、geometry.params.module\_z、materials.volume\_material\_map、source.type、source.particle、source.energy、source.position、source.direction、physics.physics\_list、output.format、output.path。
\item LLM 使用=True；回退=无；后端=llm\_slot\_frame+runtime\_semantic。
\item 第 2 轮：动作=finalize，阶段=materials，闭环状态=True。
\item 缺失字段：无。
\item LLM 使用=True；回退=无；后端=llm\_slot\_frame+runtime\_semantic。
\end{itemize}
\subsection{案例 V2C44（ev\_battery\_safety）}
\textbf{对话：}
\begin{itemize}
\item 用户：I am doing ev\_battery\_safety and need a starter simulation; ask only essential missing parameters first.
\item 系统：To set up your battery safety simulation, what basic shape should the battery geometry be, and what is its length in the X direction?
\item 用户：I need a practical medical\_radiotherapy-style setup: target as cylinder radius 49 mm half-length 52 mm made of G4\_WATER, with a beam neutron source at 2.75 MeV from (0,0,-110) mm heading (0,0,1), physics list Shielding, and hdf5 output.
\item 系统：Your simulation is ready. I've configured it as a medical radiotherapy-style setup with a cylindrical water target (radius 49 mm, half-length 52 mm). The source is a 2.75 MeV neutron beam starting at (0,0,-110) mm and directed along the positive Z-axis. The physics list is set to Shielding, and the output format is HDF5.
\end{itemize}
\textbf{内部链路（自然语言）：}
\begin{itemize}
\item 第 1 轮：动作=ask\_clarification，阶段=geometry，闭环状态=False。
\item 缺失字段：geometry.structure、geometry.params.module\_x、geometry.params.module\_y、geometry.params.module\_z、materials.volume\_material\_map、source.type、source.particle、source.energy、source.position、source.direction、physics.physics\_list、output.format、output.path。
\item LLM 使用=True；回退=无；后端=llm\_semantic\_frame+runtime\_semantic。
\item 第 2 轮：动作=finalize，阶段=materials，闭环状态=True。
\item 缺失字段：无。
\item LLM 使用=True；回退=无；后端=llm\_slot\_frame+runtime\_semantic。
\end{itemize}
\subsection{案例 V2C45（autonomous\_vehicle\_radiation）}
\textbf{对话：}
\begin{itemize}
\item 用户：??? autonomous\_vehicle\_radiation ??? Geant4 ????????????????????
\item 系统：您提到的自动驾驶车辆辐射模拟，具体需要构建哪种几何结构呢？比如，是一个简单的长方体，还是一个圆柱体？另外，这个结构的X方向尺寸是多少？
\item 用户：???? ?? 107 mm??? G4\_AIR?source.type=isotropic?particle=gamma?energy=3.50 MeV?position=(0,0,-140) mm?direction=(0,0,1)?physics=QBBC?output=xml?
\item 系统：配置已完成。
\end{itemize}
\textbf{内部链路（自然语言）：}
\begin{itemize}
\item 第 1 轮：动作=ask\_clarification，阶段=geometry，闭环状态=False。
\item 缺失字段：geometry.structure、geometry.params.module\_x、geometry.params.module\_y、geometry.params.module\_z、materials.volume\_material\_map、source.type、source.particle、source.energy、source.position、source.direction、physics.physics\_list、output.format、output.path。
\item LLM 使用=True；回退=无；后端=llm\_semantic\_frame+runtime\_semantic。
\item 第 2 轮：动作=finalize，阶段=materials，闭环状态=True。
\item 缺失字段：无。
\item LLM 使用=True；回退=无；后端=llm\_slot\_frame+runtime\_semantic。
\end{itemize}
\subsection{案例 V2C46（tokamak\_diagnostics）}
\textbf{对话：}
\begin{itemize}
\item 用户：I am doing tokamak\_diagnostics and need a starter simulation; ask only essential missing parameters first.
\item 系统：To help set up your tokamak diagnostics simulation, could you specify the basic geometry shape you'd like to start with, such as a box or cylinder?
\item 用户：I need a practical medical\_radiotherapy-style setup: target as box 180 mm x 135 mm x 60 mm made of G4\_Fe, with a point e- source at 4.25 MeV from (0,0,-170) mm heading (0,0,1), physics list FTFP\_BERT, and root output.
\item 系统：Your simulation is ready. The geometry is now a box with dimensions 180 mm x 135 mm x 60 mm, made of G4\_Fe. The source is a 4.25 MeV electron point source at (0,0,-170) mm, directed along (0,0,1). The physics list is FTFP\_BERT, and the output format is ROOT.
\end{itemize}
\textbf{内部链路（自然语言）：}
\begin{itemize}
\item 第 1 轮：动作=ask\_clarification，阶段=geometry，闭环状态=False。
\item 缺失字段：geometry.structure、geometry.params.module\_x、geometry.params.module\_y、geometry.params.module\_z、materials.volume\_material\_map、source.type、source.particle、source.energy、source.position、source.direction、physics.physics\_list、output.format、output.path。
\item LLM 使用=True；回退=无；后端=llm\_slot\_frame+runtime\_semantic。
\item 第 2 轮：动作=finalize，阶段=materials，闭环状态=True。
\item 缺失字段：无。
\item LLM 使用=True；回退=无；后端=llm\_slot\_frame+runtime\_semantic。
\end{itemize}
\subsection{案例 V2C47（lunar\_surface\_ops）}
\textbf{对话：}
\begin{itemize}
\item 用户：??? lunar\_surface\_ops ??? Geant4 ????????????????????
\item 系统：你提到的月球表面操作场景，具体需要模拟哪种几何结构呢？比如，是一个简单的长方体，还是一个圆柱体？另外，这个结构的X方向尺寸大概是多少？
\item 用户：????? ?? 73 mm ?? 88 mm??? G4\_Pb?source.type=beam?particle=proton?energy=5.00 MeV?position=(0,0,-200) mm?direction=(0,0,1)?physics=QGSP\_BERT\_HP?output=json?
\item 系统：配置已完成。
\end{itemize}
\textbf{内部链路（自然语言）：}
\begin{itemize}
\item 第 1 轮：动作=ask\_clarification，阶段=geometry，闭环状态=False。
\item 缺失字段：geometry.structure、geometry.params.module\_x、geometry.params.module\_y、geometry.params.module\_z、materials.volume\_material\_map、source.type、source.particle、source.energy、source.position、source.direction、physics.physics\_list、output.format、output.path。
\item LLM 使用=True；回退=无；后端=llm\_semantic\_frame+runtime\_semantic。
\item 第 2 轮：动作=finalize，阶段=materials，闭环状态=True。
\item 缺失字段：无。
\item LLM 使用=True；回退=无；后端=llm\_slot\_frame+runtime\_semantic。
\end{itemize}
\subsection{案例 V2C48（hospital\_physics\_qc）}
\textbf{对话：}
\begin{itemize}
\item 用户：I am doing hospital\_physics\_qc and need a starter simulation; ask only essential missing parameters first.
\item 系统：To set up your hospital physics QC simulation, what basic shape should the geometry be, and what are its dimensions along the X-axis?
\item 用户：I need a practical medical\_radiotherapy-style setup: target as sphere radius 53 mm made of G4\_W, with a isotropic neutron source at 5.75 MeV from (0,0,-230) mm heading (0,0,1), physics list Shielding, and csv output.
\item 系统：Your hospital physics QC simulation is ready.

The geometry is a medical radiotherapy-style setup with a spherical target of radius 53 mm. The target material is G4\_W (Tungsten). An isotropic neutron source with an energy of 5.75 MeV is positioned at (0,0,-230) mm, directed along (0,0,1). The physics list is set to Shielding, and the output will be in CSV format.
\end{itemize}
\textbf{内部链路（自然语言）：}
\begin{itemize}
\item 第 1 轮：动作=ask\_clarification，阶段=geometry，闭环状态=False。
\item 缺失字段：geometry.structure、geometry.params.module\_x、geometry.params.module\_y、geometry.params.module\_z、materials.volume\_material\_map、source.type、source.particle、source.energy、source.position、source.direction、physics.physics\_list、output.format、output.path。
\item LLM 使用=True；回退=无；后端=llm\_slot\_frame+runtime\_semantic。
\item 第 2 轮：动作=finalize，阶段=materials，闭环状态=True。
\item 缺失字段：无。
\item LLM 使用=True；回退=无；后端=llm\_slot\_frame+runtime\_semantic。
\end{itemize}
\subsection{案例 V2C49（synchrotron\_beamline）}
\textbf{对话：}
\begin{itemize}
\item 用户：??? synchrotron\_beamline ??? Geant4 ????????????????????
\item 系统：关于同步辐射光束线的几何结构，你希望采用哪种类型？另外，模块在X方向的尺寸（module\_x）是多少呢？
\item 用户：????? 80 mm x 180 mm x 30 mm??? G4\_Cu?source.type=point?particle=gamma?energy=6.50 MeV?position=(0,0,-80) mm?direction=(0,0,1)?physics=QBBC?output=hdf5?
\item 系统：配置已完成。
\end{itemize}
\textbf{内部链路（自然语言）：}
\begin{itemize}
\item 第 1 轮：动作=ask\_clarification，阶段=geometry，闭环状态=False。
\item 缺失字段：geometry.structure、geometry.params.module\_x、geometry.params.module\_y、geometry.params.module\_z、materials.volume\_material\_map、source.type、source.particle、source.energy、source.position、source.direction、physics.physics\_list、output.format、output.path。
\item LLM 使用=True；回退=无；后端=llm\_semantic\_frame+runtime\_semantic。
\item 第 2 轮：动作=finalize，阶段=materials，闭环状态=True。
\item 缺失字段：无。
\item LLM 使用=True；回退=无；后端=llm\_slot\_frame+runtime\_semantic。
\end{itemize}
\subsection{案例 V2C50（radiation\_safety\_drill）}
\textbf{对话：}
\begin{itemize}
\item 用户：I am doing radiation\_safety\_drill and need a starter simulation; ask only essential missing parameters first.
\item 系统：To set up your radiation safety drill simulation, what basic shape would you like to use for the shielding or detector geometry?
\item 用户：I need a practical medical\_radiotherapy-style setup: target as cylinder radius 33 mm half-length 40 mm made of G4\_Al, with a beam e- source at 7.25 MeV from (0,0,-110) mm heading (0,0,1), physics list FTFP\_BERT, and xml output.
\item 系统：Your radiation safety drill simulation is ready. I've configured a medical radiotherapy-style setup with a cylindrical aluminum target (33 mm radius, 40 mm half-length) and a 7.25 MeV electron beam starting at (0,0,-110) mm. The physics list is set to FTFP\_BERT, and output will be in XML format.
\end{itemize}
\textbf{内部链路（自然语言）：}
\begin{itemize}
\item 第 1 轮：动作=ask\_clarification，阶段=geometry，闭环状态=False。
\item 缺失字段：geometry.structure、geometry.params.module\_x、geometry.params.module\_y、geometry.params.module\_z、materials.volume\_material\_map、source.type、source.particle、source.energy、source.position、source.direction、physics.physics\_list、output.format、output.path。
\item LLM 使用=True；回退=无；后端=llm\_slot\_frame+runtime\_semantic。
\item 第 2 轮：动作=finalize，阶段=materials，闭环状态=True。
\item 缺失字段：无。
\item LLM 使用=True；回退=无；后端=llm\_slot\_frame+runtime\_semantic。
\end{itemize}
\end{document}